\documentclass[11pt]{article}
\usepackage[margin=1in]{geometry}
\usepackage{amsmath,amssymb,amsthm,mathtools}
\usepackage{enumitem}
\numberwithin{equation}{section}
\usepackage[hidelinks]{hyperref}
\hypersetup{
  pdftitle={Small Rotation Blocks and the Order-Six Stable Marriage Problem},
  pdfauthor={Myunglae Jo}
}

\newtheorem{theorem}{Theorem}[section]
\newtheorem{lemma}[theorem]{Lemma}
\newtheorem{proposition}[theorem]{Proposition}
\newtheorem{corollary}[theorem]{Corollary}
\theoremstyle{definition}
\newtheorem{definition}[theorem]{Definition}
\theoremstyle{remark}
\newtheorem{remark}[theorem]{Remark}

\newcommand{\supp}{\operatorname{supp}}
\newcommand{\Inc}{\operatorname{Inc}}
\newcommand{\down}{\mathord{\downarrow}}
\newcommand{\A}{\mathcal A}
\newcommand{\SM}{\mathcal S}

\title{Small Rotation Blocks and the Order-Six Stable Marriage Problem}
\author{Myunglae Jo}
\date{}

\begin{document}
\maketitle

\begin{abstract}
We prove that a finite poset with at most fifteen elements, width at most
three, and the following common-incomparability property has at most $48$
antichains: for every $x<y$, the elements incomparable with both $x$ and $y$
form a chain of size at most two.  An explicit fifteen-element poset attains
the bound.  We show that order-six stable-marriage rotation posets satisfy
these hypotheses by examining convex rotation blocks and their fixed partner
pools.  This gives a structural upper bound of $48$ stable matchings, matching
the lower bound already provided by Benjamin, Converse, and Krieger in 1995.
Equality forces fifteen rotations, all of size two, and every one of the $36$
man--woman pairs occurs in some stable matching.
\end{abstract}

\section{Introduction}

For a strict complete stable-marriage instance $I$, let $\SM(I)$ denote its
set of stable matchings and put
\[
        f(n)=\max_I |\SM(I)|,
\]
where the maximum is over instances with $n$ men and $n$ women.  Determining
this maximum is a classical problem; see Gusfield and Irving
\cite{GusfieldIrving1989}.  Benjamin, Converse, and Krieger
\cite{BenjaminCK1995} constructed an order-six instance with $48$ stable
matchings in 1995.  Thurber \cite{Thurber2002} subsequently developed further
lower-bound constructions.  Here we prove the matching order-six upper bound
by a local restriction on rotation posets and a sharp abstract poset
theorem.\footnote{A 2026 Biola departmental page reports the same value; at
the time of writing, we have not located a public proof and make no priority
claim.  Biola University, Department of Math and Computer Science,
``Current Students: Completed Projects,'' entry for Roman Parks (2026),
\url{https://www.biola.edu/math-and-computer-science-department/current-students},
accessed September 11, 2026.}

The use of rotation posets and their downsets is classical
\cite{IrvingLeather1986}.  Karlin, Oveis Gharan, and Weber
\cite{KarlinOGW2018} obtain an exponential bound by extracting structural
properties of these posets, and Palmer and P\'alv\"olgyi
\cite{PalmerPalvolgyi2022} give a substantially smaller exponential bound
using an entropy argument.  We use the classical rotation-poset framework,
including the downset correspondence
\cite{IrvingLeather1986,GusfieldIrving1989} and shared-agent comparability
\cite[Claim~4.6]{KarlinOGW2018}; the global rotation bound is also classical.
Our new local input concerns an order-convex block $Q$ of rotations.  If the
block uses $b$ men, a consecutive execution keeps their partners in a fixed
pool and gives $2|Q|\le b(b-1)$.  For three men, a token-return obstruction
sharpens the conclusion to $|Q|\le2$.  This differs in its hypotheses and
parameters from the support estimate in \cite[Claim~4.7]{KarlinOGW2018},
which applies to arbitrary rotation subsets and counts agents of both sexes.
For comparable rotations $\alpha<\beta$, our argument places the rotations
incomparable with both on at most $N-3$ men in an order-$N$ instance.
At $N=6$, this yields the condition called $\mathrm{CI2}$ below.

The principal counting result is that every poset with at most fifteen
elements, width at most three, and $\mathrm{CI2}$ has at most $48$
antichains.  Its proof couples the counts of two- and three-element
antichains, rather than classifying the possible posets.  We give an explicit
extremal poset and record necessary equality conditions for stable-marriage
instances.  The resulting local-to-global upper-bound proof does not use an
exhaustive classification of preference profiles.  The consecutiveness
property of an already ordered family of
maximal antichains is known \cite{DusartHC2024}; the relevant step here is
that our hypotheses order the entire family of three-element antichains
and allow the two counts to be coupled.

\section{Rotation blocks}

We use rotations and exposure in their standard sense
\cite{IrvingLeather1986,GusfieldIrving1989}; for the definition and the
elimination operation, see also \cite[Definitions~4.1 and~4.4]{KarlinOGW2018}.
We recall their properties, rather than redefine a rotation by its list
representation.  An exposed rotation is represented by a cyclic list
$((m_0,w_0),\ldots,(m_{k-1},w_{k-1}))$ of matched pairs satisfying the
standard rotation conditions, with $k\ge2$ distinct participating men and
distinct current partners.  Elimination replaces $(m_i,w_i)$ by
$(m_i,w_{i+1})$, with indices
modulo $k$, leaves nonparticipants unchanged, and yields a stable matching.
Each participating man strictly prefers $w_i$ to $w_{i+1}$.

The Irving--Leather correspondence \cite{IrvingLeather1986} identifies
stable matchings with downsets of the rotation poset.  A downset is realized
by eliminating its rotations from the man-optimal stable matching in a
linear extension of the precedence order; see \cite{GusfieldIrving1989}
and \cite[Theorem~4.5]{KarlinOGW2018}.  Rotations sharing an agent are
comparable \cite[Claim~4.6]{KarlinOGW2018}.

For a rotation $\rho$, write $\supp(\rho)$ for its set of participating men.
For $Q$ a set of rotations, put
\[
        B(Q)=\bigcup_{\rho\in Q}\supp(\rho).
\]

\begin{lemma}[Convex rotation windows]\label{lem:convex-window}
Let $Q$ be an order-convex subset of the rotation poset.  Then there are
stable matchings $M_0,M_1$ such that, starting from $M_0$, the rotations of
$Q$ can be eliminated consecutively in a linear extension of the induced
order on $Q$, producing $M_1$.
\end{lemma}

\begin{proof}
Let
\[
        I=\down Q,\qquad J=I\setminus Q.
\]
The set $I$ is a downset.  We claim that $J$ is also a downset.  Let
$y\in J$ and $x\le y$.  Since $y\in I$, there is $z\in Q$ with $y\le z$,
and hence $x\in I$.  If $x\in Q$, then
\[
        x\le y\le z,\qquad x,z\in Q,
\]
so convexity would imply $y\in Q$, a contradiction.  Hence $x\in J$.

Choose a linear extension of $J$, followed by a linear extension of the
induced order on $Q$.  This concatenation is a linear extension of $I$:
every predecessor of an element of $J$ again belongs to $J$, while a
predecessor of an element of $Q$ belongs to $J$ or to $Q$.  By the standard
realization theorem for closed sets of rotations, this gives a valid
elimination sequence for $I$.  Its $J$-prefix reaches the stable matching
corresponding to $J$, and its remaining suffix eliminates exactly the
rotations in $Q$.
\end{proof}

\begin{lemma}[Small-block capacity]\label{lem:small-block}
Let $Q$ be a finite order-convex set of rotations and put $b=|B(Q)|$.
Then every man in $B(Q)$ occurs in at most $b-1$ rotations of $Q$.  Hence
\[
        2|Q|\le b(b-1).
\]
Moreover,
\[
        b\le2\Longrightarrow |Q|\le1,
        \qquad
        b\le3\Longrightarrow |Q|\le2.
\]
\end{lemma}

\begin{proof}
Execute $Q$ as the consecutive window supplied by
Lemma~\ref{lem:convex-window}.  Let $W$ be the set of women matched to
$B(Q)$ at the beginning of this window.  Throughout the window, the men of
$B(Q)$ remain matched bijectively to exactly the same set $W$.  Indeed, at
each step the current partners of the participating subset are merely
cyclically permuted, while every other man of $B(Q)$ is unchanged.

Fix $m\in B(Q)$.  Whenever $m$ participates, his partner becomes strictly
worse; between such steps his partner is unchanged.  Thus the initial
partner of $m$ and the partners immediately following his successive
participations are all distinct elements of $W$.  Since $|W|=b$, the man
$m$ participates at most $b-1$ times.  Summing over men gives
\[
 \sum_{\rho\in Q}|\supp(\rho)|
 \le b(b-1).
\]
Every rotation has at least two participating men, proving
$2|Q|\le b(b-1)$.  The assertion for $b\le2$ follows immediately.

Suppose now that $b\le3$ and $|Q|\ge3$, and consider the first three
rotations in the window.  They contribute at least six man--rotation
incidences.  On the other hand, each of at most three men can occur at most
twice.  Equality is therefore forced: there are exactly three men, all
three rotations have support two, and every man occurs in exactly two of
the three supports.  Thus the three two-element supports form the three
edges of the triangle on these three men.  Naming $x$ as the man common to
the first and third supports and $a,b$ as the other two men, the supports
in temporal order are
\[
        \{x,a\},\qquad \{a,b\},\qquad \{x,b\}
\]
in that order.

A rotation with exactly two participating men swaps their current partners:
all nonparticipants are fixed, and the two current partner tokens must be
exchanged.  Consequently the three support-two eliminations above return
the man common to the first and third supports to his initial partner.
That man participated twice, however, and became strictly worse at each
participation.  He cannot return to his initial partner.  This contradiction
proves $|Q|\le2$.
\end{proof}

\begin{remark}\label{rem:local-global}
For the full rotation set of an order-$N$ instance,
Lemma~\ref{lem:small-block} recovers the classical bound
$|R|\le N(N-1)/2$.  The point needed here is the local version: the same
incidence argument applies to every convex rotation block, with the number
of men actually used by that block in place of $N$.  The global bound itself
is classical; see, for example, \cite{BhatnagarGR2008}.
\end{remark}

\begin{corollary}\label{cor:r15}
An order-six stable-marriage instance has at most fifteen rotations.
\end{corollary}

\begin{proof}
Apply Lemma~\ref{lem:small-block} to the entire rotation set $R$, which is
itself convex.  If $b=|B(R)|\le6$, then the proof of that lemma gives
\[
        2|R|\le b(b-1)\le6\cdot5=30.
\]
Hence $|R|\le15$.
\end{proof}

\begin{corollary}\label{cor:width3}
The rotation poset of an order-six instance has width at most three.
\end{corollary}

\begin{proof}
Rotations in an antichain cannot share a man, since rotations sharing an
agent are comparable.  Their male supports are therefore pairwise disjoint,
and each support has size at least two.  Six men can support at most three
such rotations.
\end{proof}

\section{A local structure theorem for order-six rotation posets}

For an element $x$ of a poset, let
\[
        \Inc(x)=\{z:z\parallel x\}.
\]

\begin{proposition}[Common-incomparability support budget]
\label{prop:common-budget}
Let $R$ be the rotation poset of an order-$N$ stable-marriage instance, and
let $\alpha<\beta$.  Put
\[
        Z=\Inc(\alpha)\cap\Inc(\beta).
\]
Then
\[
 |B(Z)|\le N-3,\qquad
 \operatorname{width}(Z)\le
 \left\lfloor\frac{N-3}{2}\right\rfloor,
 \qquad
 |Z|\le\binom{N-3}{2}.
\]
\end{proposition}

\begin{proof}
Let
\[
        A=[\alpha,\beta].
\]
The sets $A$ and $Z$ are convex.  Convexity of $A$ is immediate.  If
$z_1\le y\le z_2$ with $z_1,z_2\in Z$, then $y\le\alpha$ would imply
$z_1\le\alpha$, while $\alpha\le y$ would imply $\alpha\le z_2$.
Thus $y\parallel\alpha$, and the same argument with $\beta$ gives
$y\parallel\beta$.

Every $a\in A$ is incomparable with every $z\in Z$: $a\le z$ would imply
$\alpha\le z$, while $z\le a$ would imply $z\le\beta$.  Rotations sharing
a man are comparable, so
\begin{equation}\label{eq:support-disjoint}
B(A)\cap B(Z)=\varnothing.
\end{equation}
The interval $A$ contains the two distinct rotations $\alpha,\beta$.
By the two-man part of Lemma~\ref{lem:small-block}, it must therefore use
at least three men.  Equation~\eqref{eq:support-disjoint} gives
\begin{equation}\label{eq:support-budget}
|B(Z)|\le N-3.
\end{equation}

An antichain in $Z$ consists of rotations with pairwise disjoint male
supports, each of size at least two.  Hence
\[
        \operatorname{width}(Z)
        \le \left\lfloor\frac{|B(Z)|}{2}\right\rfloor
        \le \left\lfloor\frac{N-3}{2}\right\rfloor.
\]
Finally $Z$ is convex, so Lemma~\ref{lem:small-block} gives
\[
        2|Z|\le |B(Z)|(|B(Z)|-1),
\]
which together with \eqref{eq:support-budget} yields the stated cardinality bound.
\end{proof}

\begin{theorem}[Common incomparability]\label{thm:ci2}
Let $R$ be the rotation poset of an order-six stable-marriage instance.
For every $\alpha<\beta$ in $R$, the set
\[
        Z(\alpha,\beta)=\Inc(\alpha)\cap\Inc(\beta)
\]
is a chain and has cardinality at most two.
\end{theorem}

\begin{proof}
Proposition~\ref{prop:common-budget} gives $|B(Z)|\le3$ and
$\operatorname{width}(Z)\le1$, so $Z$ is a chain.  Since $Z$ is convex,
the three-man part of Lemma~\ref{lem:small-block} gives $|Z|\le2$.
\end{proof}

This explains why order six is the natural endpoint of the present local
argument.  For $Z=\Inc(\alpha)\cap\Inc(\beta)$,
Proposition~\ref{prop:common-budget} gives $|B(Z)|\le N-3$,
$\operatorname{width}(Z)\le\lfloor(N-3)/2\rfloor$, and
$|Z|\le\binom{N-3}{2}$.
At $N=6$, the bounds $|B(Z)|\le3$ and $\operatorname{width}(Z)\le1$
force a chain on at most three participating men, and the exceptional
three-man capacity sharpens the cardinality bound to $|Z|\le2$.
At $N=7$, the same general budget gives only $|B(Z)|\le4$,
$\operatorname{width}(Z)\le2$, and $|Z|\le6$.
Thus this mechanism forces the $\mathrm{CI2}$ rigidity used below at order
six, but the same budget no longer forces it at order seven.

Theorem~\ref{thm:ci2} motivates the following purely order-theoretic
condition.

\begin{definition}\label{def:ci2}
A finite poset $P$ has property $\mathrm{CI2}$ if, whenever $x<y$,
the common incomparability set
\[
        \Inc(x)\cap\Inc(y)
\]
is a chain of cardinality at most two.
\end{definition}

\begin{proposition}[Forbidden-pattern form]\label{prop:forbidden}
A finite poset satisfies $\mathrm{CI2}$ if and only if it has no induced
subposet isomorphic to $2+1+1$ or to $2+3$, where the summands are mutually
incomparable chains.
\end{proposition}

\begin{proof}
If $\mathrm{CI2}$ fails for $x<y$, then either
$\Inc(x)\cap\Inc(y)$ contains two incomparable elements, which together
with $x<y$ induce $2+1+1$, or it contains a three-element chain, which
together with $x<y$ induces $2+3$.  The converse is immediate from either
forbidden configuration.
\end{proof}

\begin{remark}\label{rem:graph-ci2}
Equivalently, the incomparability graph of a $\mathrm{CI2}$ poset has no
induced diamond and no induced $K_{2,3}$; antichains of the poset are
precisely cliques of its incomparability graph.
\end{remark}

\section{A finite-poset bound}

Let $P$ be a finite poset of width at most three, let $n=|P|$, and let
$A_j$ denote the number of $j$-element antichains.  Since every antichain
has size at most three,
\begin{equation}\label{eq:antichain-decomposition}
|\A(P)|=1+n+A_2+A_3.
\end{equation}

We first record the two bounds needed for the generic case.

\begin{lemma}[Triple structure]\label{lem:triple}
Assume $P$ has width at most three and satisfies $\mathrm{CI2}$.  If
$A_3>0$, then
\[
        A_3\le \left\lfloor\frac{n-1}{2}\right\rfloor.
\]
\end{lemma}

\begin{proof}
Let $T$ be a three-element antichain and $x\notin T$.  The point $x$ is
comparable with at least one member of $T$, else $T\cup\{x\}$ would be a
four-element antichain.  In fact it is comparable with at least two.
Otherwise, if $x$ were comparable with only $t\in T$, the other two
members of $T$ would be incomparable elements of
$\Inc(x)\cap\Inc(t)$, contradicting $\mathrm{CI2}$.  Moreover $x$ cannot
lie above one member of $T$ and below another, since those two members
would then be comparable.  Thus every $x\notin T$ lies below at least two
members of $T$, or above at least two.

Write these alternatives as $x\in L(T)$ and $x\in U(T)$.  If
$l\in L(T)$ and $u\in U(T)$, the two corresponding two-element subsets of
$T$ intersect, so $l<t<u$ for some $t\in T$; hence
\begin{equation}\label{eq:lower-upper-separation}
L(T)<U(T).
\end{equation}

Now let $S,T$ be three-element antichains.  We claim that their generated
ideals are nested.  Suppose first that some $u\in S\setminus T$ lies in
$U(T)$.  Then no element of $S\setminus T$ can lie in $L(T)$, because
\eqref{eq:lower-upper-separation} would make it comparable with $u$.  Thus every element of
$S\setminus T$ lies in $U(T)$.

We show that every $t\in T$ lies in $\down S$.  This is immediate if
$t\in S$.  Otherwise suppose that $t\not\le s$ for every $s\in S$.
For $s\in S\cap T$ we have $t\parallel s$.  For
$s\in S\setminus T$, the point $s$ lies above at least two members of
$T$.  Hence $t>s$ is impossible: among those two members choose
$t'\ne t$; then $t'<s<t$, contradicting that $T$ is an antichain.
Our assumption $t\not\le s$ therefore forces $t\parallel s$ for every
$s\in S$.  Then $S\cup\{t\}$ is a four-element antichain, contradicting
the width bound.  Hence $t\le s$ for some $s\in S$, and so
$\down T\subseteq\down S$.

If no element of $S\setminus T$ lies in $U(T)$, then every such element
lies in $L(T)$ and hence below a member of $T$.  Thus
$S\subseteq\down T$ and $\down S\subseteq\down T$.  This proves nesting.

Distinct three-element antichains generate distinct ideals, since the
maximal elements of $\down T$ are exactly the members of the antichain
$T$.  Also two distinct triples share at most one point: if $S\ne T$
shared two points, a point of $S\setminus T$ would be incomparable with
two members of $T$, contrary to the first paragraph.

Order all three-element antichains as
\[
        T_1,\ldots,T_t,\qquad
        \down T_1\subsetneq\cdots\subsetneq\down T_t,
        \quad t=A_3.
\]
We shall use the following no-re-entry property, the usual consecutiveness
property for ordered maximal antichains; see
\cite[Proposition~2.4 in the preprint version]{DusartHC2024}.
It applies here because a three-element antichain is maximal in a
width-three poset, and the maximal-antichain order agrees with inclusion
of generated ideals.  We include the short proof for completeness.
If $i<j<k$ and $x\in T_i\cap T_k$, then
$x\in T_j$.  Indeed $x\in\down T_j$, so $x\le y$ for some $y\in T_j$;
also $y\in\down T_k$, so $y\le z$ for some $z\in T_k$.  Since
$x,z\in T_k$ and $T_k$ is an antichain, $x=y=z$.

It follows that each $T_{i+1}$ meets
$T_1\cup\cdots\cup T_i$ in at most one point.  The first triple contributes
three points and every later triple contributes at least two new points.
Hence
\[
        n\ge 3+2(t-1)=2t+1.
\]
\end{proof}

\begin{lemma}[Pair bound]\label{lem:pairs}
Under the assumptions of Lemma~\ref{lem:triple}, if $A_3>0$, then
\[
        A_2\le 2n-3.
\]
\end{lemma}

\begin{proof}
By Dilworth's theorem, partition $P$ into three chains
$C_1,C_2,C_3$.  Because a three-element antichain exists, each chain is
nonempty.  Consider one chain
\[
        C=\{c_1<\cdots<c_a\}.
\]
For every $x\notin C$, the set of elements of $C$ incomparable with $x$
is an interval: if $c_i\le c_j\le c_k$ and $x\parallel c_i,c_k$, then
$x\parallel c_j$.

Indeed, if the incomparable interval of an outside point $x$ contains
$r\ge1$ elements of $C$, then $x$ contributes $r$ to the left-hand side
and $1+(r-1)$ to the right-hand side of
\begin{equation}\label{eq:chain-interval-count}
\sum_{c\in C}|\Inc(c)|
 =
 N_C+\sum_{i=1}^{a-1}
 |\Inc(c_i)\cap\Inc(c_{i+1})|,
\end{equation}
where $N_C$ is the number of points outside $C$ incomparable with at least
one point of $C$.  Here $N_C\le n-a$, while every adjacent intersection
in \eqref{eq:chain-interval-count} has size at most two by $\mathrm{CI2}$.  Therefore
\[
        \sum_{c\in C}|\Inc(c)|\le n+a-2.
\]
Summing over the three chains counts every incomparable pair exactly twice.
If $|C_1|+|C_2|+|C_3|=n$, then
\[
        2A_2\le (n+|C_1|-2)+(n+|C_2|-2)+(n+|C_3|-2)
        =4n-6,
\]
which proves the claim.
\end{proof}

\begin{corollary}[Linear antichain growth]\label{cor:linear-growth}
If $P$ has width three and satisfies $\mathrm{CI2}$, then
\begin{equation}\label{eq:linear-bound}
|\A(P)|
        \le
        3n-2+\left\lfloor\frac{n-1}{2}\right\rfloor.
\end{equation}
In particular, the number of antichains is linear in $n$.
\end{corollary}

\begin{proof}
Width three means $A_3>0$.  Substitute the bounds of
Lemmas~\ref{lem:triple} and \ref{lem:pairs} into
$|\A(P)|=1+n+A_2+A_3$.
\end{proof}

\begin{remark}
The linear behavior is a genuine restriction.  The disjoint union of three
chains of balanced lengths has width three and
\[
        (|C_1|+1)(|C_2|+1)(|C_3|+1)=\Theta(n^3)
\]
antichains.  Thus $\mathrm{CI2}$ collapses the possible antichain count from
cubic to linear growth within width three.
\end{remark}

For $n\le14$, Corollary~\ref{cor:linear-growth} already gives a bound
strictly below $48$.  The only remaining issue is the endpoint $n=15$,
where the generic linear estimate gives $50$.  We need one small tradeoff.

\begin{lemma}[The fifteen-point refinement]\label{lem:n15}
Let $P$ satisfy the hypotheses above with $n=15$.  If $A_3\ge6$, then
\[
        |\A(P)|\le46.
\]
\end{lemma}

\begin{proof}
Order the three-element antichains
$T_1,\ldots,T_t$ as in Lemma~\ref{lem:triple}, put
\[
        U_i=T_1\cup\cdots\cup T_i,\qquad U=U_t,
\]
and let $m$ be the number of indices $i<t$ for which
$T_{i+1}\cap U_i=\varnothing$.  By no re-entry,
$T_{i+1}\cap U_i\subseteq T_i$, and distinct triples meet in at most one
point.  Therefore
\begin{equation}\label{eq:triple-union-card}
|U|=2t+1+m.
\end{equation}

Let $e(U)$ denote the number of incomparable pairs with both endpoints in
$U$.

\emph{Claim 1.}  At a transition $T_i\to T_{i+1}$, the number of new cross
incomparable pairs is zero if $T_i\cap T_{i+1}\ne\varnothing$, and at most
three if $T_i\cap T_{i+1}=\varnothing$.

To localize the possible new cross pairs, let
\[
 x\in U_i\setminus T_i,\qquad
 y\in T_{i+1}\setminus U_i,
\]
choose $j<i$ with $x\in T_j$.  Since $\down T_j\subseteq\down T_i$
and $x\notin T_i$, some member of $T_i$ lies strictly above $x$;
Lemma~\ref{lem:triple} therefore places $x$ in $L(T_i)$.
The point $y\in T_{i+1}\setminus U_i$ cannot lie below a member
$z$ of $T_i$: the ideal inclusion would give $z\le w$ for some
$w\in T_{i+1}$, contradicting that $T_{i+1}$ is an antichain.
Thus Lemma~\ref{lem:triple} places $y$ on the upper side of $T_i$,
so $y\in U(T_i)$.  By \eqref{eq:lower-upper-separation}, $x<y$.
Thus every new incomparability joining a new vertex to the old union
has its old endpoint in $T_i$.  In what follows, a cross pair excludes
pairs already internal to the new triple; in particular a shared vertex
is not an old-only endpoint.

If $T_i$ and $T_{i+1}$ meet in one point $s$, let
$x\in T_i\setminus\{s\}$ and
$y\in T_{i+1}\setminus\{s\}$.  The point $x$ lies in $L(T_{i+1})$ and is
incomparable with $s$, so it must lie below both other members of
$T_{i+1}$; in particular $x<y$.  Hence a nonempty transition creates no
cross incomparability.

If the two consecutive triples are disjoint, their bipartite
incomparability graph is a matching: each point outside a triple is
incomparable with at most one member of that triple.  Hence an empty
transition creates at most three cross incomparabilities.  The new triple
itself contributes three incomparable pairs in either case.  The accounting
is therefore
\[
\begin{array}{c|c|c}
|T_i\cap T_{i+1}| & \text{new internal pairs} & \text{new cross pairs}\\
\hline
1 & 3 & 0\\
0 & 3 & \le3
\end{array}
\]
Induction over the transitions gives
\begin{equation}\label{eq:union-pair-bound}
e(U)\le3t+3m.
\end{equation}

\emph{Claim 2.}  Every $x\notin U$ has at most two incomparable neighbors
in $U$:
\begin{equation}\label{eq:outsider-neighbor-bound}
|\Inc(x)\cap U|\le2.
\end{equation}

For every triple
$T_i$, the first paragraph of
Lemma~\ref{lem:triple} puts $x$ either in $L(T_i)$ or in $U(T_i)$.
These alternatives are monotone along the ideal order: $x$ cannot be below
an earlier triple $S$ and above a later triple $T$.  Indeed choose
$s\in S$ with $x<s$.  Since $\down S\subseteq\down T$, choose
$u\in T$ with $s\le u$.  If $x$ were above at least two members of $T$,
one of them, say $v$, would differ from $u$, and
$v<x<s\le u$ would contradict the antichain property of $T$.

Thus the triples above which $x$ lies form an initial segment and the
triples below which $x$ lies form a final segment.  If both are nonempty,
let $T_c$ be the last former triple and $T_{c+1}$ the first latter triple.
Every vertex in an earlier triple but outside $T_c$ lies below $T_c$ and
therefore below $x$ by \eqref{eq:lower-upper-separation}; dually every later vertex outside
$T_{c+1}$ lies above $x$.  Hence all incomparable neighbors of $x$ in
$U$ lie in $T_c\cup T_{c+1}$.  The first paragraph of
Lemma~\ref{lem:triple} gives at most one such neighbor in each triple.
The endpoint cases are handled by the same localization using $T_1$ or
$T_t$.  This proves Claim~2.

Put $q=15-|U|$.  Combining
\eqref{eq:union-pair-bound} and \eqref{eq:outsider-neighbor-bound} with
\eqref{eq:triple-union-card} and the trivial bound on pairs among the $q$
outside points gives
\begin{equation}\label{eq:pair-tradeoff}
A_2\le3t+3m+2q+\binom q2,
        \qquad
        q=14-2t-m.
\end{equation}
Lemma~\ref{lem:triple} gives $t\le7$.  If $t=7$, then $m=q=0$ and
$A_2\le21$.  If $t=6$, then $m+q=2$; the possibilities
$(m,q)=(0,2),(1,1),(2,0)$ give respectively
$A_2\le23,23,24$.  In all cases $t\ge6$,
\[
        A_2+A_3\le30,
\]
and \eqref{eq:antichain-decomposition} yields $|\A(P)|\le46$.
\end{proof}

We also need the elementary width-two endpoint.

\begin{lemma}[Width two]\label{lem:width2}
If $P$ has width at most two, satisfies $\mathrm{CI2}$, and $n\ge2$, then
\[
        |\A(P)|
        \le 2n+\left\lfloor\frac n2\right\rfloor-1.
\]
In particular, $|\A(P)|\le36$ when $n\le15$.
\end{lemma}

\begin{proof}
Partition $P$ into two nonempty chains $C,D$ of sizes $a\le b$,
$a+b=n$.  (If $P$ is itself a chain, split it into two nonempty
subchains.)  As in \eqref{eq:chain-interval-count},
\[
        A_2=\sum_{c\in C}|\Inc(c)|
        \le b+2(a-1).
\]
Since $a\le\lfloor n/2\rfloor$,
\[
        |\A(P)|=1+n+A_2
        \le 2n+a-1
        \le 2n+\left\lfloor\frac n2\right\rfloor-1.
\]
\end{proof}

\begin{remark}
The endpoint $n=15$ is the only place where the separate estimates on
$A_2$ and $A_3$ are insufficient.  Lemma~\ref{lem:n15} uses the ordered
family of three-element antichains to couple the two counts.
\end{remark}

\begin{theorem}[The finite-poset theorem]\label{thm:poset48}
Let $P$ be a finite poset satisfying
\[
        |P|\le15,\qquad \operatorname{width}(P)\le3,
        \qquad \mathrm{CI2}.
\]
Then
\[
        |\A(P)|\le48.
\]
\end{theorem}

\begin{proof}
The cases $n\le1$ are immediate.  If $A_3=0$, the width is at most two
and Lemma~\ref{lem:width2} applies.

Assume $A_3>0$.  If $n\le14$, then
Corollary~\ref{cor:linear-growth} gives $|\A(P)|\le46$.
Now let $n=15$.  If $A_3\ge6$, use Lemma~\ref{lem:n15}.  If
$1\le A_3\le5$, Lemma~\ref{lem:pairs} gives $A_2\le27$, and hence
\[
        |\A(P)|=16+A_2+A_3\le16+27+5=48.
\]
\end{proof}

\begin{proposition}[An explicit extremal poset]\label{prop:explicit-extremal}
The constant $48$ in Theorem~\ref{thm:poset48} is attained by
$P_* = \{1,\ldots,5\}\times\{1,2,3\}$ with strict order
\begin{equation}\label{eq:explicit-order}
(i,a)<(j,b)\quad\Longleftrightarrow\quad
j\ge i+2\ \text{or}\ (j=i+1\ \text{and}\ a\ne b).
\end{equation}
\end{proposition}

\begin{proof}
The relation is irreflexive.  In two successive strict comparisons the
level increases by at least two in total, which proves transitivity.
Thus adjoining equality gives a partial order.
Within each level all three points are incomparable.  Between consecutive
levels exactly the three same-coordinate pairs are incomparable; points
at more distant levels are comparable.  A mixed-level three-element
antichain is impossible: the singleton in one level would have to share
its second coordinate with both points in the other level.  It follows that
the only three-element antichains are the five levels, and none has more
than three elements.  Hence
\begin{equation}\label{eq:explicit-counts}
       A_2=5\binom32+4\cdot3=27,\qquad A_3=5,
       \qquad |\A(P_*)|=1+15+27+5=48.
\end{equation}
To check $\mathrm{CI2}$, take $(i,a)<(j,b)$.  If $j=i+1$, then
$a\ne b$ and the common incomparability set is
$\{(i,b),(i+1,a)\}$, a two-element chain.  If $j=i+2$, that set is
$\{(i+1,a)\}$ when $a=b$, and empty otherwise.  If $j\ge i+3$, it is
empty.  This checks every comparable pair.
\end{proof}

\section{The order-six maximum}

\begin{theorem}\label{thm:main}
The maximum number of stable matchings in a strict complete stable-marriage
instance with six men and six women is
\[
        f(6)=48.
\]
\end{theorem}

\begin{proof}
Let $R$ be the rotation poset of an arbitrary order-six instance.
Corollaries~\ref{cor:r15} and \ref{cor:width3}, together with
Theorem~\ref{thm:ci2}, give
\[
        |R|\le15,\qquad
        \operatorname{width}(R)\le3,\qquad
        R\text{ satisfies }\mathrm{CI2}.
\]
By Theorem~\ref{thm:poset48}, $R$ has at most $48$ antichains.

For a finite poset, downsets are in bijection with antichains via maximal
elements and generated ideals.  By the Irving--Leather rotation theorem,
stable matchings of the instance are in bijection with downsets of $R$.
Thus the instance has at most $48$ stable matchings.

The order-six construction of Benjamin, Converse, and Krieger already
has $48$ stable matchings \cite{BenjaminCK1995}.  Hence $f(6)=48$.
\end{proof}

\begin{corollary}[Necessary conditions for equality]\label{cor:equality}
If an order-six instance has $48$ stable matchings, then its rotation poset
$R$ satisfies
\begin{equation}\label{eq:equality-counts}
        |R|=15,\qquad A_2=27,\qquad A_3=5.
\end{equation}
Every rotation has exactly two participating men, every man belongs to
exactly five rotations, and every man--woman pair occurs in a stable matching.
\end{corollary}

\begin{proof}
The preceding bounds give at most $46$ antichains when $|R|\le14$,
and at most $36$ at width at most two.  Thus $|R|=15$ and $A_3>0$.
Lemma~\ref{lem:n15} excludes $A_3\ge6$ at equality.  Since
$A_2\le27$ and $A_3\le5$, the identity $48=16+A_2+A_3$ forces
$A_2=27$ and $A_3=5$.

Apply the participation bound to a complete elimination sequence.  For
each of the six men let $d_m$ be his total number of participations, with
$d_m=0$ for any man in no rotation.  Then $d_m\le5$, and
\begin{equation}\label{eq:equality-incidences}
30=2|R|\le\sum_{\rho\in R}|\supp(\rho)|
       =\sum_m d_m\le6\cdot5=30.
\end{equation}
Every inequality is an equality.  Consequently each support has size two
and each $d_m$ equals five.  A man sees his initial partner and five
subsequent, strictly worse partners along the sequence.  These six distinct
partners exhaust the six women, and every intermediate matching is stable.
Thus every man--woman pair occurs in some stable matching.
\end{proof}

\begin{remark}[The classical example and $P_*$]\label{rem:ds3-pstar}
The Benjamin--Converse--Krieger $DS_3$ instance also realizes the abstract
sharp example of Proposition~\ref{prop:explicit-extremal}.  A direct finite
reconstruction of its stable-matching lattice gives $15$ join-irreducibles
whose poset is isomorphic to $P_*$; an explicit isomorphism and reproducible
certificate are given in the Supplementary Material.  This observation is
not used in the upper-bound proof.
\end{remark}

\section{Verification note}

The new finite combinatorial core has accompanying Lean verification.
The three formal layers cover the fixed-token participation and small-block
bounds, the support-poset implication from explicit capacity hypotheses to
width at most three and $\mathrm{CI2}$, and the finite-poset bound of
Theorem~\ref{thm:poset48}, including its fifteen-element endpoint.
The formalization is a verification aid, not a logical dependency of the
manuscript proof.
The connection from actual rotations to the token model and the standard
stable-matching--downset correspondence is handled in the manuscript rather
than formalized in Lean.
The Supplementary Material records the precise formalization scope and
reproducibility details.
The accompanying Lean sources and reproducibility artifacts are available at
\url{https://github.com/insammael/stable-marriage-f6}.

\paragraph{Declarations.} Generative-AI tools were used as aids for proof exploration, Lean
formalization, literature searching, code preparation, and manuscript
drafting and revision.  The author assumes full responsibility for all
content.

\begin{thebibliography}{99}

\bibitem{BenjaminCK1995}
A. T. Benjamin, C. Converse, and H. A. Krieger,
How do I marry thee? Let me count the ways!,
\emph{Discrete Applied Mathematics} 59 (1995), 285--292.
doi:10.1016/0166-218X(95)80006-P.

\bibitem{IrvingLeather1986}
R. W. Irving and P. Leather,
The complexity of counting stable marriages,
\emph{SIAM Journal on Computing} 15 (1986), 655--667.
doi:10.1137/0215048.

\bibitem{DusartHC2024}
J. Dusart, M. Habib, and D. G. Corneil,
Maximal cliques lattices structures for cocomparability graphs with
algorithmic applications,
\emph{Order} 41 (2024), 99--133.
doi:10.1007/s11083-023-09641-x.
Preprint version: \emph{Maximal cliques structure for cocomparability
graphs and applications}, arXiv:1611.02002 (2016).

\bibitem{PalmerPalvolgyi2022}
C. Palmer and D. P\'alv\"olgyi,
At most $3.55^n$ stable matchings,
in \emph{Proceedings of FOCS 2021}, 217--227, IEEE, 2022.
doi:10.1109/FOCS52979.2021.00029.

\bibitem{GusfieldIrving1989}
D. Gusfield and R. W. Irving,
\emph{The Stable Marriage Problem: Structure and Algorithms},
MIT Press, 1989.

\bibitem{KarlinOGW2018}
A. R. Karlin, S. Oveis Gharan, and R. Weber,
A simply exponential upper bound on the maximum number of stable matchings,
in \emph{Proceedings of the 50th Annual ACM SIGACT Symposium on Theory of
Computing (STOC 2018)}, 920--925, 2018.
doi:10.1145/3188745.3188848.

\bibitem{BhatnagarGR2008}
N. Bhatnagar, S. Greenberg, and D. Randall,
Sampling stable marriages: why spouse-swapping won't work,
in \emph{Proceedings of the Nineteenth Annual ACM--SIAM Symposium on
Discrete Algorithms (SODA 2008)}, 1223--1232, 2008.
doi:10.1145/1347082.1347215.

\bibitem{Thurber2002}
E. G. Thurber,
Concerning the maximum number of stable matchings in the stable marriage
problem,
\emph{Discrete Mathematics} 248 (2002), 195--219.
doi:10.1016/S0012-365X(01)00194-7.

\end{thebibliography}

\end{document}
